% ── TC Metric Definitions — Method Section ────────────────────────────────────
% Place inside your Methods chapter/section.
% Requires: amsmath

\subsection{Temporal Consistency Metrics}

To quantify jitter in robot action trajectories, seven temporal consistency (TC)
metrics were evaluated. All metrics are defined such that a higher value
corresponds to greater jitter, and are computed per joint then averaged across
all $J = 6$ joints. Let $a_j[t]$ denote the commanded joint angle for joint $j$
at timestep $t$, with $T$ total timesteps and sampling period $\Delta t =
1/30$~s.

\subsubsection{Mean Absolute Consecutive Difference (MACD)}

\begin{equation}
    \text{MACD} = \frac{1}{(T-1)\,J}
    \sum_{t=1}^{T-1} \sum_{j=1}^{J} \bigl|a_j[t] - a_j[t-1]\bigr|
    \label{eq:macd}
\end{equation}

MACD measures the average step-to-step change in commanded joint position. It is
simple and parameter-free, directly reflecting how much each action command
differs from the previous one.

\subsubsection{Normalised Jerk}

Jerk is the third derivative of position with respect to time and is a
well-established measure of motion smoothness in robotics and motor
control~\cite{flash1985coordination}. It is approximated here using finite
differences:

\begin{equation}
    \text{Jerk} = \frac{1}{J} \sum_{j=1}^{J}
    \frac{1}{T-3} \sum_{t} \left|
        \frac{\Delta^3 a_j[t]}{\Delta t^3}
    \right| \cdot \frac{1}{\mathrm{range}(a_j)}
    \label{eq:jerk}
\end{equation}

where $\Delta^3$ denotes the third-order finite difference and
$\mathrm{range}(a_j) = \max_t a_j[t] - \min_t a_j[t]$ normalises by the
signal range to make the metric scale-independent across joints.

\subsubsection{Direction Change Count (DCC)}

\begin{equation}
    \text{DCC} = \frac{1}{J} \sum_{j=1}^{J}
    \frac{1}{T-2} \sum_{t=2}^{T-1}
    \mathbf{1}\bigl[\,\mathrm{sgn}(\Delta a_j[t]) \neq \mathrm{sgn}(\Delta a_j[t-1])\,\bigr]
    \label{eq:dcc}
\end{equation}

DCC counts the fraction of timesteps at which the direction of motion reverses
for each joint. This captures oscillatory back-and-forth behaviour that may be
missed by frequency-domain metrics when the oscillation amplitude is small.

\subsubsection{Spectral Entropy (SE)}

Spectral entropy applies the Shannon entropy~\cite{shannon1948mathematical} to
the normalised power spectrum of the action signal:

\begin{equation}
    \text{SE} = -\frac{1}{\log_2(N/2+1)}
    \sum_{f} p(f)\log_2 p(f),
    \qquad
    p(f) = \frac{|F(f)|^2}{\sum_f |F(f)|^2}
    \label{eq:se}
\end{equation}

where $F(f)$ is the discrete Fourier transform of the mean-subtracted signal
and the prefactor normalises SE to $[0, 1]$. A smooth signal concentrates
energy at low frequencies (low SE); a jittery signal spreads energy across the
spectrum (high SE).

\subsubsection{Weighted Mean Frequency (Sm)}

\begin{equation}
    \text{Sm} = \frac{2}{N\,f_s} \sum_{f} |F(f)| \cdot f
    \label{eq:sm}
\end{equation}

Sm is the amplitude-weighted mean frequency of the action signal, where $f_s =
30$~Hz is the sampling rate and the factor $2/N$ normalises for the FFT length.
Higher Sm indicates that more signal energy resides at high frequencies,
corresponding to jittery motion. Sm is used as the primary TC metric throughout
training and evaluation in this work.

\subsubsection{Signal Energy Ratio (SER)}

\begin{equation}
    \text{SER} = \frac{\sum_{f \leq f_{\text{th}}} |F(f)|^2}{\sum_{f} |F(f)|^2}
    \label{eq:ser}
\end{equation}

SER measures the fraction of signal energy contained below a threshold
frequency $f_{\text{th}} = 1$~Hz, chosen to separate intentional low-frequency
robot motion from high-frequency jitter. Higher raw SER indicates a smoother
signal; in the comparative evaluation (Table~\ref{tab:tc_real}), $1-\text{SER}$
is reported to maintain the convention that higher values correspond to greater
jitter.

\subsubsection{Autocorrelation Decay Rate (ADR)}

ADR fits an exponential decay to the normalised autocorrelation function:

\begin{equation}
    R(\tau) \approx e^{-\lambda\tau}
    \implies
    \lambda = -\frac{\mathrm{d}}{\mathrm{d}\tau}\log R(\tau)
    \label{eq:adr}
\end{equation}

The decay constant $\lambda$ (s$^{-1}$) is estimated by least-squares regression
on $\log R(\tau)$ over lags $\tau \in [0,\, \min(T/2,\, 50)]\,\Delta t$. A
smooth trajectory remains self-correlated over long lags (low $\lambda$); a
jittery trajectory loses correlation quickly (high $\lambda$). ADR was found to
be non-monotonic on synthetic data and was excluded from the primary evaluation.

% ── Bibliography entries (add to your .bib file) ──────────────────────────────
%
% @article{flash1985coordination,
%   author  = {Flash, Tamar and Hogan, Neville},
%   title   = {The coordination of arm movements: an experimentally confirmed
%              mathematical model},
%   journal = {Journal of Neuroscience},
%   volume  = {5},
%   number  = {7},
%   pages   = {1688--1703},
%   year    = {1985}
% }
%
% @article{shannon1948mathematical,
%   author  = {Shannon, Claude E.},
%   title   = {A mathematical theory of communication},
%   journal = {Bell System Technical Journal},
%   volume  = {27},
%   number  = {3},
%   pages   = {379--423},
%   year    = {1948}
% }
